\lyseis{7.2}{Απόλυτη τιμή, αντίθετοι, σύγκριση}
\sol{1} α) $7$ \quad β) $7$ \quad γ) $0$ \quad δ) $2{,}5$ \quad ε) $\frac{3}{4}$ \quad στ) $0{,}01$.
\sol{2} $-5$, $+8$, $0$, $-2{,}5$, $+\frac{1}{3}$, $-0{,}6$.
\sol{3} α) $4$ \quad β) $-3$ \quad γ) $-2$ \quad δ) $-6$.
\sol{4} α) $5$ ή $-5$ \quad β) $0$ \quad γ) $\frac{1}{2}$ ή $-\frac{1}{2}$ \quad δ) αδύνατη.
\sol{5} α) $-3>-8$ \quad β) $2>-9$ \quad γ) $-0{,}5<-0{,}05$ \quad δ) $-\frac{7}{2}=-3{,}5<-3$ \quad ε) $0>-0{,}001$ \quad στ) $-4<4$.
\sol{6} $-7<-3<-0{,}5<0<2{,}5<4<5$.
\sol{7} Μεγαλύτερος ο $-1{,}02$ και μικρότερος ο $-1{,}22$.
\sol{8} α) Σωστό. β) Λάθος: $|-5|=5$. γ) Σωστό. δ) Λάθος: $-9<-2$. ε) Σωστό. στ) Σωστό.
\sol{9} $-10<-5<-3<-1<7$: Φλώρινα, Μέτσοβο, Ιωάννινα, Θεσσαλονίκη, Αθήνα.
\sol{10} Το βιβλιοπωλείο ($|{-2}|=2$ μονάδες $=200$ μ.) είναι πιο κοντά από τον φούρνο ($300$ μ.). Η απόσταση μεταξύ τους είναι από το $-2$ στο $+3$: $5$ μονάδες $=500$ μ.
\sol{11} $-11\,034<-430<0<2\,917$.
\sol{12} Χρωστάει περισσότερο η Ζωή ($|{-15}|=15$ €). Το μεγαλύτερο ποσό έχει η Ρήνη ($+12$).
\sol{13} Το $\Lambda$ έχει τετμημένη $+7$. Η απόσταση $K\Lambda$ είναι $14$ μονάδες ($7$ αριστερά και $7$ δεξιά από το $0$).
\sol{14} Στον άξονα το $-10$ είναι πολύ πιο αριστερά από το $-3$, άρα μικρότερο. Στο θερμόμετρο $10$ βαθμούς κάτω από το μηδέν κάνει περισσότερο κρύο από $3$ βαθμούς κάτω. Το $10>3$ ισχύει για τις απόλυτες τιμές, όχι για τους αρνητικούς.
\sol{15} Το $\alpha$ και το $-\alpha$ έχουν την ίδια απόσταση από το $0$.
\sol{16} Το $|\alpha|=\alpha$ ισχύει για $\alpha\ge0$ και το $|\alpha|=-\alpha$ για $\alpha\le0$.
\sol{17} Για παράδειγμα $x=3$, $y=-3$. Ακέραιοι με απόλυτη τιμή μικρότερη από $2$: $-1,\,0,\,1$, δηλαδή $3$.
\sol{18} Για παράδειγμα $-1{,}5$ και $-0{,}25$, $-0{,}22$. Υπάρχουν άπειροι.
\sol{19} Ελεύθερη απάντηση.
